\section{Introduction}

\subsection{Background}
Automation is contributing to a decline in skilled technical labor, while current robotic systems remain insufficient to fully replace human operators. As a result, demand for skilled workers persists. However, acquiring such skills remains difficult, particularly due to the need for prolonged training and the development of muscle memory.

Learning in technical domains consists of both conceptual understanding and physical skill acquisition. The latter, often requiring repeated practice and embodied interaction, presents a significant barrier to entry.

\subsection{Previous Studies}
Existing research explores XR as a medium for immersive training in domains such as medical procedures and flight simulation. Additionally, AI techniques have been applied to control systems and training optimization.

However, current drone simulators (e.g., entertainment-oriented platforms) primarily target experienced users, featuring steep learning curves and limited onboarding support. They also rely heavily on specialized hardware, creating accessibility barriers for beginners.

\subsection{Limitations of Existing Approaches}
Prior work largely emphasizes concept transfer and high-level training outcomes, with less focus on early-stage skill acquisition and muscle memory formation.

Furthermore, existing simulators lack:
\begin{itemize}
    \item Guided interaction for beginners
    \item Accessible control interfaces
    \item Integration of AI for real-time assistance
\end{itemize}

\subsection{Research Gap}
There is limited research on using XR as a tool specifically for lowering the barrier of **muscle memory acquisition** through guided, interactive, and iterative training.

\subsection{Problem Statement}
How can XR be designed to support accessible and effective muscle memory training for complex physical skills?

\subsection{Goal}
To develop an XR-based training system that reduces the learning curve of FPV drone piloting through AI-assisted co-training and interaction design.

\subsection{Research Questions}
\begin{itemize}
    \item How can XR environments facilitate muscle memory development?
    \item How can AI assist users in real-time during physical skill training?
    \item What interaction design supports gradual skill acquisition?
\end{itemize}

\subsection{Methods Overview}
This study adopts a design-based research approach with iterative prototyping. It focuses on the development of an XR FPV drone simulator integrating physics simulation, AI-assisted co-training (PID and Reinforcement Learning), and interaction design for guided practice.